\documentclass[12pt]{article}
\usepackage{fancyhdr}
\pagestyle{fancy}
\fancyhf{}
\fancyfoot[C]{\thepage}
\fancyfoot[L]{\scriptsize CC BY-NC-SA 4.0 \textbullet\ Leo Liang}
\renewcommand{\headrulewidth}{0pt}
\usepackage[margin=1in]{geometry}
\usepackage{amsmath, amssymb}
\usepackage{array, booktabs}
\usepackage{pgffor}

\parindent=0pt
\parskip=1em
\pagestyle{fancy}

% ─────────────────────────────────────────────────────────────────────────────
\begin{document}

\begin{center}
  {\Large\bfseries The Mathematician Who Counted the Uncountable}\\[4pt]
  Georg Cantor (1845--1918)\\[6pt]
  \textit{Week 2 History Spotlight}
\end{center}
\hrule\bigskip

\noindent\textbf{The Setup}

In 1874, a German mathematician named Georg Cantor published a paper that asked
one of the strangest questions in the history of mathematics: \textit{are all
infinities the same size?}

Everyone assumed the answer was yes. Infinity is infinity. It's the biggest thing
there is. What would it even mean for one infinity to be bigger than another?

Cantor proved they were wrong. Some infinities are strictly larger than others.
And to do it, he invented set theory --- the same language you are learning today.

\bigskip
\noindent\textbf{Where He Came From}

Cantor was born in Saint Petersburg, Russia, to a Danish merchant father and a
Russian mother. The family moved to Germany when he was eleven, and he spent most
of his life there. From an early age he was a gifted student, but unremarkable in
the eyes of the mathematical establishment --- a detail that would matter a lot later.

He got his doctorate from the University of Berlin in 1867, writing a thesis on
number theory. Respectable work. Nothing that hinted at what was coming.

\bigskip
\noindent\textbf{The Idea That Changed Everything}

Cantor started thinking about a deceptively simple question: how do you compare
the ``sizes'' of two infinite sets?

For finite sets, it's easy --- you count. But what does it mean to count an infinite
collection? You can't just say ``the natural numbers are bigger than the even numbers''
because both go on forever.

Cantor's insight was to use \textbf{pairing}. Two sets have the same size if you
can perfectly match their elements, one-to-one, with nothing left over.
In the language you're learning: they have the same size if there's a \textbf{bijection}
between them.

Using this idea, he showed something wild: the natural numbers $\mathbb{N}$ and
the even natural numbers $\{2, 4, 6, 8, \ldots\}$ are actually the \textbf{same
size}. You can pair them:

\[
  1 \leftrightarrow 2,\quad
  2 \leftrightarrow 4,\quad
  3 \leftrightarrow 6,\quad
  4 \leftrightarrow 8,\quad \ldots
\]

Every natural number gets a partner. No natural number is left out. No even number
is doubled up. It's a perfect \textbf{bijection}. So the two sets are the same ``size,''
even though one seems smaller.

This was already strange. But Cantor wasn't done.

\bigskip
\noindent\textbf{The Diagonal Argument}

Cantor then asked: what about the real numbers $\mathbb{R}$? Could we also pair
every real number with a natural number?

He proved we cannot. No matter how cleverly you try to list the real numbers between
0 and 1, he could always find a real number you missed. His method --- now called
\textbf{Cantor's diagonal argument} --- is still one of the most beautiful proofs
in all of mathematics. It goes roughly like this:

\textbf{Suppose} you claim you have a complete list of all real numbers between 0 and 1:

\medskip
\begin{center}
\begin{tabular}{@{} c l @{}}
  1st number: & $0.\mathbf{3}14159\ldots$ \\
  2nd number: & $0.7\mathbf{1}8281\ldots$ \\
  3rd number: & $0.57\mathbf{7}215\ldots$ \\
  4th number: & $0.141\mathbf{5}92\ldots$ \\
  \vdots & \vdots
\end{tabular}
\end{center}

\medskip
Now build a new number: look at the \textbf{diagonal} digits (bolded above:
$3, 1, 7, 5, \ldots$) and change each one. If it's 3, write 4. If it's 1, write 2. If it's 9, write 0.
And so on. The number you just built is different from the 1st entry in the 1st decimal place,
different from the 2nd in the 2nd place, different from the 3rd in the 3rd place
--- different from \textit{every} entry on the list, in the one spot that matters.

So your ``complete'' list was not complete. There's always a real number missing.
Cantor concluded that the real numbers are \textbf{strictly more numerous}
than the natural numbers. A larger kind of infinity.

He called the size of $\mathbb{N}$ the first infinite cardinal, $\aleph_0$
(``aleph-null''), and the size of $\mathbb{R}$ something larger: $2^{\aleph_0}$.
He spent years trying to prove there was no infinity ``in between'' these two.
That question --- the \textbf{Continuum Hypothesis} --- was never resolved in his lifetime.
(It is still, in a precise sense, unresolvable. But that's a story for another day.)

\bigskip
\noindent\textbf{The Reception}

The mathematical world did not applaud.

Leopold Kronecker --- one of the most powerful mathematicians in Germany,
and Cantor's former professor --- despised the work. He called Cantor a
``corrupter of youth.'' He publicly ridiculed him, blocked him from getting
professorships in Berlin, and wrote letters campaigning against his publications.

His objection was philosophical: he believed mathematics should only deal with
things you can explicitly construct, step by step. Infinite sets, in his view,
were nonsense. ``God made the integers,'' Kronecker famously said,
``all else is the work of man.''

Cantor was devastated. He suffered what his doctors called ``nervous exhaustion''
--- what we would today recognize as severe depression. He was hospitalized multiple
times over the years in a sanatorium in Halle.

Between episodes, he kept working. He published. He corresponded with mathematicians
across Europe who understood and supported him. He wrote long, impassioned letters
defending his ideas. He also developed a keen interest in Shakespeare scholarship
--- he became convinced, seriously, that Francis Bacon had secretly written
Shakespeare's plays, and spent years on that project too. (He was wrong about this one.)

\bigskip
\noindent\textbf{The End, and the Legacy}

Cantor died in 1918 in the sanatorium where he had been living.
He had spent the final years of his life asking to be allowed to go home.
He died before he could.

Within a generation, the mathematical world had completely reversed its verdict.
David Hilbert --- the most influential mathematician of the early 20th century ---
called Cantor's work ``the finest product of mathematical genius and one of the
supreme achievements of purely intellectual human activity.''

Today, set theory is the foundation of essentially all of modern mathematics.
When you write $f \colon A \to B$ or talk about domain and codomain,
you are using Cantor's language. Every theorem in analysis, algebra, and topology
is built on the idea that mathematical objects can be organized into sets and
studied by their relationships.

Cantor did not live to see that vindication.

\bigskip
\noindent\textbf{One Thing to Think About}

Cantor showed that there are different sizes of infinity by finding a function ---
a bijection (or proving one can't exist). The same tool he used to compare sets
is the same tool we use to understand functions.

When you prove that a function is injective or surjective, you are asking Cantor's
question in miniature: can these two sets be perfectly paired?

\bigskip
\hrule\bigskip
\noindent \textit{Discussion question}: Kronecker said you should only
do mathematics with things you can explicitly build, step by step.
Cantor said infinite sets are real mathematical objects even if you can't list them all.
Who do you think is right? Does it matter?

\bigskip

Write a response in your words. 

\end{document}
